\documentclass[11pt,a4paper]{article}
\usepackage[T1]{fontenc}
\usepackage{lmodern}
\usepackage[margin=25mm]{geometry}
\usepackage{amsmath,amssymb,amsthm,microtype}
\usepackage[dvipsnames]{xcolor}
\usepackage[colorlinks=true,linkcolor=MidnightBlue,urlcolor=MidnightBlue,citecolor=MidnightBlue]{hyperref}
\pagestyle{plain}
\setlength{\emergencystretch}{2em}
\newtheorem{theorem}{Theorem}[section]
\newtheorem{lemma}[theorem]{Lemma}
\newtheorem{corollary}[theorem]{Corollary}
\theoremstyle{remark}\newtheorem{remark}[theorem]{Remark}
\newcommand{\dd}{\,\mathrm d}
\newcommand{\e}{\mathrm e}
\newcommand{\D}{\mathcal D}
\newcommand{\Rea}{\operatorname{Re}}
\hypersetup{pdftitle={Importing arithmetic windows into the prime-power cone},
pdfauthor={Robert Polak},
pdfsubject={A short-window criterion for exclusion of zeta zeros and its consequence for the prime-power cone}}
\title{Importing arithmetic windows\\into the prime-power cone}
\author{Ing.\ Robert Polak\thanks{Correspondence:
\href{mailto:robopol@gmail.com}{robopol@gmail.com}.}\\[0.4em]
\small Mathematical derivations and manuscript preparation with GPT-6 Pro (OpenAI)\\
\small Mathematical review and editing with GPT ASTRA 6.2 (OpenAI)}
\date{21 November 2026 \quad (version 1.0)}
\begin{document}\maketitle
\begin{abstract}
Fix \(0\le s<1/2\). We show that a one-sided bound
\(F(x)\le C_s x^{1/2+s}\) for the integrated von Mangoldt residual on
arbitrarily late intervals \([Z,Z\log Z]\) excludes every zeta zero
with real part greater than \(1/2+s\). The proof uses the positive
kernel of the completed zeta function to isolate a single zero,
including its multiplicity. The resulting oscillation has exponential
size on each sufficiently late interval of logarithmic time of length
\(\log(T/2)\) centred at \(T\).
We apply this criterion to an arithmetic window estimate extracted
from the Alp\"oge--Claude development. Taking that estimate as input
gives RH and the strict determinant bound from the earlier prime-power
cone criterion. We identify the Lean declarations supplying the
arithmetic input and give the analytic argument here in ordinary
mathematical prose.
\end{abstract}

\section{The arithmetic interface to import}
Define
\begin{equation}
 F(x)=\sum_{n\le x}\Lambda(n)(x/n-1)
                 -x\log x+(1+\gamma_E)x.\label{eq:F}
\end{equation}
This is exactly \texttt{Formalism.Gz} in the released source
\cite{Source}, \texttt{Formalism/StatementOC.lean}.

The application to the Alp\"oge--Claude development uses the following
arithmetic estimate.
\medskip
\noindent\textbf{Arithmetic input (A).}
There exist fixed \(E>0\) and \(\varepsilon_*>0\) such that for every
\(0<\varepsilon<\varepsilon_*\) and every \(Z_0\) there exists
\(Z\ge\max(Z_0,2)\) for which
\begin{equation}
 F(x)\le x^{1/2+E\varepsilon}
       \qquad(Z\le x\le Z^{3/2}).\label{eq:external}
\end{equation}
The constants \(E\) and \(\varepsilon_*\) are chosen before
\(\varepsilon\). The subsequent construction parameters and
thresholds may depend on \(\varepsilon\). No uniformity as
\(\varepsilon\downarrow0\) is required.
\medskip

In the source this is obtained by combining the arithmetic
comparison, the existence of suitable families and their windows,
and the principal-character specialization. It is not presented
there under this exact single theorem name. Section
\ref{sec:audit} records the extraction boundary. In particular,
the assertion is stronger than the conditional statement of
\texttt{core\_gm} alone: the existence hypotheses of that statement
must be supplied by the preceding construction.

The concrete source uses \(e_D=2\) and \(e_L=0\), which gives
\(E=e_D+e_L+7=9\). Extracting those exact witnesses would require
a strengthened wrapper. For the existing public interface, it is
enough to take \(E=\max(1,e_D+e_L+7)\) and use \(x\ge1\).
The analytic proof below needs only this fixed-positive-\(E\) form.

The cutoff \(\varepsilon_*\) incorporates internal positive
thresholds. The explicit restriction \(\varepsilon<1/52\) visible
in window selection is not, by itself, the complete permitted range
of all later steps. Only the existence of a positive cutoff matters
in the argument below.

\paragraph{Relation to the existing paper.}
Section~9 of the companion paper \cite{Polak} explains that
the arithmetic window bound, together with the external windowed
Landau theorem, yields RH and hence the cone conclusion. The
contribution of the present note is the direct xi-kernel proof of
the window-to-spectrum step and the explicit import boundary.

\section{The retained growth factor}
Use the residual and notation
\begin{align}
 B(t)&=\sum_\rho a_\rho\e^{\lambda_\rho t},&
 a_\rho&=\frac1{\rho(1-\rho)},&
 \lambda_\rho&=\rho-\frac12,\label{eq:B}\\
 C&=B(0)=2+\gamma_E-\log(4\pi),&
 \Phi(t)&=C-B(t).\nonumber
\end{align}
Zeros are counted with multiplicity. The classical zero count gives
\(S:=\sum_\rho|a_\rho|<\infty\). The function \(B\) is real and
even, and
\begin{equation}
 |B(v)|\le S\e^{|v|/2}\quad(v\in\mathbb R).\label{eq:Bbound}
\end{equation}
The exact prime-side formula of \cite{Polak} implies
\begin{equation}
 B(\log x)=F(x)/\sqrt{x}+O(x^{-1/2})\quad(x\to\infty).
 \label{eq:BF}
\end{equation}
The correction is unconditional and uniform for large real \(x\).

\begin{lemma}[Exponential oscillation from one zero]\label{lem:exp}
Suppose \(\rho_*\) is one nontrivial zero with
\(\alpha:=\Rea\rho_*-1/2\ge0\). Then there are constants
\(b_*>0\) and \(T_*>0\), depending on this zero, such that every
\(T\ge T_*\) admits points \(t_+,t_-\) in
\begin{equation}
 [T-R_T,T+R_T],\qquad R_T=\tfrac12\log(T/2),\label{eq:R}
\end{equation}
satisfying
\begin{equation}
 B(t_+)\ge b_*\e^{\alpha T},\qquad
 B(t_-)\le-b_*\e^{\alpha T}.\label{eq:exposc}
\end{equation}
No assumption about a rightmost zero or simple zeros is needed.
\end{lemma}
\begin{proof}
We use the classical positive even xi kernel \cite{Newman}. For
\(u\ge0\), write
\[
 K(u)=\sum_{n\ge1}
   (4\pi^2n^4\e^{9u/2}-6\pi n^2\e^{5u/2})
                 \e^{-\pi n^2\e^{2u}},
\]
extended evenly. Its transform and tails satisfy
\begin{align}
 M(z):=\int_{\mathbb R}K(u)\e^{zu}\dd u&=\xi(1/2+z),\nonumber\\
 c_0\e^{9|u|/2-\pi\e^{2|u|}}
 &\le K(u)\le C_0\e^{9|u|/2-\pi\e^{2|u|}}.\label{eq:K}
\end{align}
Here \(\xi(w)=\tfrac12 w(w-1)\pi^{-w/2}\Gamma(w/2)\zeta(w)\).
To check the normalisation, set
\(f(u)=\e^{u/2}\sum_{n\ge1}\exp(-\pi n^2\e^{2u})\) on the real line.
The Jacobi theta transformation gives
\(f(-u)-f(u)=(\e^{u/2}-\e^{-u/2})/2\). Thus \(f''-f/4\) is even,
and termwise differentiation for \(u\ge0\) identifies it with \(K\).
Writing \(w=1/2+z\), Mellin integration gives
\(\int_{\mathbb R}f(u)\e^{zu}\dd u
=\tfrac12\pi^{-w/2}\Gamma(w/2)\zeta(w)\) for \(\Rea w>1\).
Two integrations by parts supply the factor \(z^2-1/4=w(w-1)\).
The kernel's double-exponential decay then extends the transform
identity to every \(z\).
For \(u\ge0\) each summand of \(K\) is positive. Keeping the first
summand gives \(c_0=4\pi^2-6\pi>0\), and dropping the negative terms
inside the brackets gives
\(C_0=4\pi^2\sum_{n\ge1}n^4\e^{-\pi(n^2-1)}\).

Put \(\lambda=\rho_*-1/2\), and let \(r\) be the exact
multiplicity of \(\rho_*\). Define
\begin{align}
 h_r(u)
 &=\frac1{(r-1)!}\int_0^\infty
       v^{r-1}\e^{\lambda v}K(u+v)\dd v\nonumber\\
 &=\frac{(-1)^r}{(r-1)!}\int_0^\infty
       v^{r-1}\e^{-\lambda v}K(u-v)\dd v.\label{eq:h}
\end{align}
The equality follows by splitting
\(\int_{\mathbb R}(w-u)^{r-1}\e^{\lambda w}K(w)\dd w=0\)
at \(w=u\). This integral vanishes because
\(M^{(j)}(\lambda)=0\) for \(0\le j<r\).

The first formula bounds the right tail and the second the left.
Using \(\e^{2v}-1\ge2v\) and \(|\Rea\lambda|\le1/2\),
\[
 \frac{|h_r(u)|}{K(u)}
 \le \frac{C_0/c_0}{(2\pi\e^{2|u|}-5)^r}\le L<\infty.
\]
Here \(2\pi-5>0\). Consequently the transform of \(h_r\)
is entire. Absolute Fubini for \(\Rea z>\alpha\), followed by
analytic continuation, gives
\begin{equation}
 H_r(z)=\int_{\mathbb R}h_r(u)\e^{zu}\dd u
       =\frac{M(z)}{(z-\lambda)^r},\qquad
 H_r(\lambda)=\frac{M^{(r)}(\lambda)}{r!}\ne0.\label{eq:H}
\end{equation}

Equations \eqref{eq:Bbound} and \eqref{eq:K}, together with
\(|h_r|\le LK\), justify termwise integration of the full zero
series, including the region \(T+u<0\). Thus
\begin{align}
 \int K(u)B(T+u)\dd u&=0,\label{eq:zeroavg}\\
 \int h_r(u)B(T+u)\dd u
   &=A_*\e^{\lambda T},&
 A_*&=r a_{\rho_*}\frac{M^{(r)}(\lambda)}{r!}\ne0.
 \label{eq:nonzeroavg}
\end{align}
The factor \(r\) comes from multiplicity in the zero sum.
Let \(d_*=|A_*|>0\), and put \(B_+=\max(B,0)\),
\(B_-=\max(-B,0)\). Crucially, retain the full magnitude in
\eqref{eq:nonzeroavg}:
\begin{equation}
 \int K B_+(T+u)\dd u
 =\int K B_-(T+u)\dd u
 \ge\frac{d_*}{2L}\e^{\alpha T}.\label{eq:massgrowth}
\end{equation}
For a zero on the critical line this gives a fixed lower bound.
For a zero to its right, \(\alpha>0\) and both integrals grow at
least exponentially with \(T\).

For \(R\ge0\), the tails satisfy
\begin{align*}
 \int_{|u|>R}K(u)|B(T+u)|\dd u
 &\ll\e^{T/2}\int_R^\infty\e^{5v-\pi\e^{2v}}\dd v\\
 &\ll\e^{T/2+3R-\pi\e^{2R}}.
\end{align*}
With \(R=R_T\), this is
\(O(T^{3/2}\e^{-(\pi-1)T/2})=o(1)\), hence also
\(o(\e^{\alpha T})\). Both central signed integrals are
therefore at least \(d_*\e^{\alpha T}/(4L)\) for large \(T\).
Set \(M_0=\int K=\xi(1/2)>0\) and
\(b_*=d_*/(4LM_0)\). The corresponding maximum of \(B_+\)
and of \(B_-\) on the compact interval is at least
\(b_*\e^{\alpha T}\), proving \eqref{eq:exposc}.
\end{proof}

\section{Sparse arithmetic windows exclude zeros}
\begin{theorem}[A window-to-spectrum implication]\label{thm:window}
Fix \(s\in[0,1/2)\) and \(C_s>0\). Suppose there are
arbitrarily large \(Z\) such that
\begin{equation}
 F(x)\le C_s x^{1/2+s}\qquad(Z\le x\le Z\log Z).
 \label{eq:logwindow}
\end{equation}
Then every nontrivial zero satisfies \(\Rea\rho\le1/2+s\).
In particular, the conclusion holds if \eqref{eq:logwindow}
is known on the longer windows \([Z,Z^{1+\kappa}]\), for any
fixed \(\kappa>0\).
\end{theorem}
\begin{proof}
Suppose instead that a zero has
\(\alpha=\Rea\rho-1/2>s\). Apply Lemma \ref{lem:exp} to
this one zero. Choose one of the arbitrarily late windows, put
\(U=\log Z\), and set
\[
 T=U+\tfrac12\log U,\qquad R_T=\tfrac12\log(T/2).
\]
For large \(U\), \(T\le2U\), so \(R_T\le\tfrac12\log U\).
It follows that
\[
 [T-R_T,T+R_T]\subseteq[U,U+\log U].
\]
The point \(t_+\) in Lemma \ref{lem:exp} therefore lies within
the logarithmic image of the arithmetic window. By
\eqref{eq:BF} and \eqref{eq:logwindow},
\[
 b_*\e^{\alpha T}\le B(t_+)
       \le C_s'\e^{s t_+}
       \le C_s'\e^{sT}(T/2)^{s/2}.
\]
Thus
\(\e^{(\alpha-s)T}\le(C_s'/b_*)(T/2)^{s/2}\), which is
impossible along the arbitrarily late windows. This excludes
every such zero individually. Finally,
\(Z\log Z\le Z^{1+\kappa}\) for all sufficiently large \(Z\),
so the longer-window assertion follows.
\end{proof}

\begin{corollary}[Arithmetic input (A) implies RH]\label{cor:RH}
If the imported statement \eqref{eq:external} holds for every
sufficiently small fixed \(\varepsilon>0\), then RH holds.
\end{corollary}
\begin{proof}
For every such \(\varepsilon\) with \(E\varepsilon<1/2\),
Theorem \ref{thm:window} gives
\(\Rea\rho\le1/2+E\varepsilon\) for every zero. Equivalently,
given a hypothetical zero with displacement \(\alpha>0\),
first choose one fixed
\[
 0<\varepsilon<\min(\varepsilon_*,\alpha/(2E)).
\]
Fix all parameters for this \(\varepsilon\), and only then take
an arbitrarily late window. Its exponent satisfies
\(E\varepsilon<\alpha\), contradicting Theorem \ref{thm:window}.
No rightmost zero or uniform threshold in \(\varepsilon\) has
been assumed. The functional equation completes RH.
\end{proof}

The proof of Corollary~\ref{cor:RH} uses the external arithmetic
construction only through (A). Lemma~\ref{lem:exp} supplies the
analytic zero-exclusion step.

Conversely, if all nontrivial zeros satisfy \(\Rea\rho\le1/2+s\),
then \(|B(t)|\le S\e^{st}\) for \(t\ge0\) by absolute convergence.
Equation~\eqref{eq:BF} gives \(F(x)\ll x^{1/2+s}\) for all large \(x\).
In particular, RH gives \(F(x)\le A\sqrt{x}\) for some \(A>0\).
For each fixed \(\varepsilon>0\), \(x^\varepsilon\ge A\) eventually,
so (A) follows with \(E=1\). Thus (A) and RH are equivalent for
this particular residual. The arithmetic content of the application
lies in obtaining (A) from the external construction.

\section{The cone conclusion and the earlier obstructions}
Use Polak's definitions
\[
 k(t)=\e^{t/2}-\frac{\e^{-t/2}}{\e^{2t}-1}-\frac C4,
 \qquad
 \D_\pm(t)=\Phi(t)(2k(t)+\Phi(t)/4)-\Phi'(t\pm)^2.
\]
Once Corollary \ref{cor:RH} places the zeros on the line,
the absolutely convergent spectral representation gives
\begin{equation}
 \Phi(t)=2\sum_{\gamma>0}
      \frac{1-\cos(\gamma t)}{\gamma^2+1/4},
 \qquad 0\le\Phi(t)\le2C.\label{eq:phibound}
\end{equation}
The conditional cone theorem of \cite{Polak} now gives
\(\D_+(\log q)>0\) at all sufficiently large prime powers
\(q\). We give the derivation of the sharper asymptotic bound
recorded in the working note \cite{PreviousNote}.
\begin{corollary}\label{cor:margin}
Under RH, and hence under arithmetic input (A), for either one-sided
derivative one has
\[
 \D_\pm(t)\ge(1/32-o(1))t^2.
\]
\end{corollary}
\begin{proof}
Put \(X=\e^{t/2}\) and \(h=t/(4X)\). The cosine identity and the
expansion at the origin in Sections~3 and~4 of \cite{Polak} give
\[
 4\Phi(t)\ge2\Phi(h)-\bigl[\Phi(t+h)+\Phi(t-h)-2\Phi(t)\bigr],
 \qquad 2\Phi(h)=h\log(1/h)+O(h).
\]
On \([t-h,t+h]\) the smooth density of \(\Phi''\) is
\(k(v)+\Phi(v)/4\le(1+o(1))X\), and every atom is negative.
Integration against the triangular weight \((h-|v-t|)_+\) bounds
the second difference above by \((1+o(1))Xh^2\). Since
\(\log(1/h)=t/2-\log(t/4)\), this gives
\(\Phi(t)\ge(1/64-o(1))t^2/X\) for every sufficiently large real \(t\).

For either slope \(p=\Phi'(t\pm)\), the RH Chebyshev estimate in
\cite[Section~4]{Polak} gives \(|p|=O(t^2)\). Take \(u=-p/k(t)\), so
\(|u|=O(t^2/X)\). The upper Taylor estimate from the same section,
valid across downward jumps for either choice of initial slope, gives
\[
 \Phi(t+u)\le\Phi(t)-\frac{p^2}{2k(t)}
                       +O(u^2+X|u|^3).
\]
The error is \(o(t^2/X)\). The lower bound just proved, applied at
\(t+u=t+o(1)\), gives
\(\Phi(t+u)\ge(1/64-o(1))t^2/X\). Multiplication by
\(2k(t)=(2+o(1))X\) and addition of \(\Phi(t)^2/4\) yield the result.
\end{proof}

\paragraph{A short direct bound for the weak target.}
Even without the strict cone reserve, \eqref{eq:phibound} and
the prime-side measure identity already supply the desired weak
floor. Put \(X=\e^{t/2}\) and let \(p\) be either one-sided slope
of \(\Phi\). On \(|v-t|\le1\), its smooth second-derivative
density is at most \(2X\) for large \(t\), and all its derivative
jumps are negative. Integral Taylor expansion in the direction
opposite to \(p\) gives, for \(0<h<1\),
\[
 0\le\Phi(t)-|p|h+Xh^2\le2C-|p|h+Xh^2.
\]
The inequality works for either initial side of a downward jump.
Taking \(h=\sqrt{2C/X}\) gives \(p^2\le8CX\). Since \(k(t)>0\)
eventually,
\begin{equation}
 \D_\pm(t)\ge-8C\e^{t/2}.\label{eq:weak}
\end{equation}
Consequently the weak sampled target
\(d_m=\D_+(\log p_m)/\sqrt{p_m}\ge-K\) holds at every
sufficiently large index, with \(K=8C\). With the strict cone
theorem, \(K=0\) works as well. The waiting time for this weak
target is then zero.

\paragraph{Relation to the short-interval problem.}
The local transfer in \cite{Polak} retains the factor
\(x^{E\varepsilon}\) inside an arithmetic window.
Here \(\varepsilon\) is fixed for one hypothetical zero. The forbidden
excursion occurs in every sufficiently late interval of the length in
Lemma~\ref{lem:exp}, so an arithmetic window beyond the fixed thresholds
suffices. The spectral conclusion then applies at every prescribed sample.
The joint error estimate for four shifted prime sums on a prescribed
prefix remains a separate question. For example, the even profile
\(B_*(t)=|t|\) has the same large-positive-time growth as the model
\(F_*(x)=\sqrt{x}\log x\) in \cite{Polak}. Its average against any
translate of \(K\) is strictly positive, whereas the zeta residual
satisfies \eqref{eq:zeroavg}.

\section{Source extraction and verification boundary}\label{sec:audit}
This section records a static reading of the released source at
commit \texttt{40e56bab} of \cite{Source}. It distinguishes the
arithmetic theorem we import from the analytic statements that
follow it in that source. A standalone Lean declaration for
\eqref{eq:external} has not been written or compiled here.

\subsection{The comparison and its arithmetic hypotheses}
\texttt{UpperTheory.Comparison\textsubscript{2}.core\_gm}
in \href{https://github.com/Deicyde/riemann-hypothesis/blob/40e56bab/src/Development/AutomorphicArithmetic/PrimitiveCosets/CoreGm.lean}{\texttt{CoreGm.lean}}
has the following essential quantifier order: after fixing the
level, \(e_D,e_L\), the comparison constants and
\(\varepsilon>0\), it supplies a threshold \(x_1\); then it
bounds \(\Lambda.\mathrm{Gm}(x)\) for all \(x\ge x_1\) and all
arithmetic frames satisfying \texttt{EAtH} and \texttt{CoreAtH}.
Its exponent is
\[
 \frac12+(e_D+e_L+7)\varepsilon.
\]
The threshold is chosen before the arithmetic witnesses, so those
witnesses may vary with \(x\) across one full window.

This is a conditional comparison theorem, not a construction of
its hypotheses. Its underlying \texttt{ch\_top} calls
\texttt{each\_unit\_bound\_statement}, whose conclusion has the
form \(ASK\le Main+Err\). The input includes divisor, height,
slope, loss, counting and parameter bounds.
\texttt{PChain.errX5at\_le} absorbs the error at sufficiently
large \(x\), including the requirements
\[
 \log x\le x^\varepsilon,\qquad
 200(1+K)^2(1+c_{HH})(1+e_K)\le x^\varepsilon.
\]
The scale condition \(|a(t)/b-x|\le1\) is retained.
The heavy arithmetic content lies in realizing these hypotheses.

\subsection{Families, complete windows and fixed parameters}
\href{https://github.com/leanprover-community/riemann-hypothesis/blob/40e56bab/src/Development/AutomorphicArithmetic/PrimitiveCosets/WindowInequalityOfArtinInputs.lean}{\texttt{windowInequality\_of\_artinInputs}}
and
\href{https://github.com/leanprover-community/riemann-hypothesis/blob/40e56bab/src/Development/AutomorphicArithmetic/PrimitiveCosets/LevelGmLeOnFamilyScales.lean}{\texttt{level\_gm\_le\_on\_family\_scales}}
supply the property \texttt{ChebyshevInequalityOnWindows} defined in
\href{https://github.com/Deicyde/riemann-hypothesis/blob/40e56bab/src/Formalism/SharpRestriction/RamificationFrobenius/TotientGaloisClasses/TopWin.lean}{\texttt{TopWin.lean}}.
This supplies the same bound on a family's windows above a
single threshold. The definition \texttt{HasWindows} in
\href{https://github.com/Deicyde/riemann-hypothesis/blob/40e56bab/src/Formalism/SharpRestriction/RamificationFrobenius/TotientGaloisClasses/Window.lean}{\texttt{Window.lean}}
means
\[
 \forall y\ \exists Y\ge y\ \forall x\in[Y,Y^{\ell}],
 \quad x\text{ belongs to the family's window set}.
\]
It therefore supplies full real intervals and absorbs any fixed
large-\(x\) threshold by requesting a later \(Y\).

\texttt{FamilyClassConnected}, defined in
\href{https://github.com/Deicyde/riemann-hypothesis/blob/40e56bab/src/Development/AutomorphicArithmetic/CharacterBoxes/DivinNOfHypV.lean}{\texttt{DivinNOfHypV.lean}},
chooses fixed \(e_D,e_L\) and a positive epsilon cutoff
\emph{before} choosing epsilon-dependent tied families. The
current top-level path is
\texttt{zero\_free\_all\_of\_green\_v\_win} in
\href{https://github.com/Deicyde/riemann-hypothesis/blob/40e56bab/src/Development/AutomorphicArithmetic/ThetaCorrespondence/ThetaPairLifts/ThirdPlaceSeries/TopPlug.lean}{\texttt{ThirdPlaceSeries/TopPlug.lean}}.
Before its zero-free call it constructs \texttt{TopG} from
\texttt{forall\_goal\_of\_shape\_of\_hyps} and
\texttt{shape\_closed\_link\_n\_of\_green}; \texttt{supplyChi}
supplies \texttt{HostExW} for \(\ell=3/2\).

The concrete proof in
\href{https://github.com/Deicyde/riemann-hypothesis/blob/40e56bab/src/Development/AutomorphicArithmetic/ThetaCorrespondence/ThetaPairLifts/ThirdPlaceSeries/TopShape.lean}{\texttt{ThirdPlaceSeries/TopShape.lean}}
ends with \texttt{topDesWin} applied to parameters
\(2,3,2,0\); the last two are \(e_D=2,e_L=0\).
\texttt{topDesWin} returns those witnesses unchanged, with a
positive minimum of internal epsilon cutoffs. Its public return
type hides their exact numerical values. This explains both the
concrete coefficient \(9\) and the safer public-interface
form with an unspecified fixed \(E>0\).

\subsection{Specialization and the precise stopping point}
Apply the resulting \texttt{TopG} at the trivial primitive
character, using \texttt{hypV\_of\_cls} and
\texttt{eFamC\_of\_cls}. These are the family-producing steps in
\href{https://github.com/Deicyde/riemann-hypothesis/blob/40e56bab/src/Development/AutomorphicArithmetic/CharacterBoxes/ZeroFreeAllOfTopG.lean}{\texttt{ZeroFreeAllOfTopG.lean}}.
For \texttt{Level.cyc PChar.one}, the subgroup generated by the
trivial character consists of that character alone. Its retained
primitive character set is \(\{1\}\), and erasing \texttt{PChar.one}
leaves an empty sum. The definitions of
\href{https://github.com/leanprover-community/riemann-hypothesis/blob/40e56bab/src/Formalism/SharpRestriction/RamificationFrobenius/TotientGaloisClasses/ExplicitFormulaErrorOfHostConditions.lean}{\texttt{Level.cyc}, \texttt{Level.Nt} and \texttt{Level.Gm}}
therefore reduce the aggregate to \texttt{Formalism.Gz},
which is exactly \eqref{eq:F}. Combining this with the preceding
comparison and absorbing its threshold gives arithmetic input (A).

This extraction stops \emph{before} \texttt{formB},
\texttt{zeroFree\_of\_rpow\_window}, \texttt{rh\_of\_every},
or the final RH theorem. Theorem \ref{thm:window} supplies our
own analytic conversion of the arithmetic input instead.

\subsection{Where the existence premises are supplied}
The starting point is
\href{https://github.com/leanprover-community/riemann-hypothesis/blob/40e56bab/src/Development/AutomorphicArithmetic/CharacterBoxes/MainConclusion.lean}{\texttt{MainConclusion.lean}}.
Its proof uses a row from the theta-lift construction and the existence
result \texttt{recTExists} to assemble the literature record. The same
row supplies \texttt{VWin} through \texttt{vwin\_of\_row}.
The row itself is constructed in
\href{https://github.com/leanprover-community/riemann-hypothesis/blob/40e56bab/src/Development/AutomorphicArithmetic/Foundations/GaloisRepresentations/IntegralStructures/RowSrc.lean}{\texttt{RowSrc.lean}}
by applying \texttt{galoisRepresentations\_of\_thetaLift}.
The next step supplies the theta integrals, local matching data and
geometric covolume bounds. These are used to construct the
hypothesis record passed to \texttt{TopPlug} and \texttt{TopShape}.

These proof terms supply the premises of \texttt{core\_gm}.
The interfaces traced here require no explicit assumption of RH.
Our source review covers this extraction, including the order in
which its witnesses are chosen.

\subsection{Reported checks versus checks performed here}
Andrew Sutherland's MIT report \cite{MIT}, updated
21 November 2026 at 08:45 UTC, reports that all 175,187 modules
compile, all have passed Lean kernel replay, and the five headline
theorems depend only on Lean's three standard axioms. It reports
acceptance of the full export by nanoda and con-leche, agreement of
the theorem statements under con-leche's comparison, and successful
module checks by lean4lean. The end-to-end \texttt{lake comparator}
run and Lean-kernel replay of the full export were still running
at that update. These are reported third-party checks.

In this note the extraction of (A) was checked by source inspection.
The new analytic argument was checked on paper. A machine-checked
version would also require formalising that argument and compiling
the specialised statement (A).

\section{Conclusion and attribution}
For the classical zeta function, the positive xi kernel forces
oscillations of size \(\e^{\alpha T}\) from each zero with real part
\(1/2+\alpha\). Their localisation is strong enough to use bounds
on \([Z,Z\log Z]\). In the application, the longer arithmetic
windows are supplied by the Alp\"oge--Claude development. The
determinant and its differential identities come from \cite{Polak}.
Version 1.2 already described how external arithmetic windows and
a separate windowed Landau theorem would imply the cone result.
The proof here supplies the analytic implication directly and gives
the shorter sufficient window length.

\begin{thebibliography}{9}
\bibitem{Polak}
R.~Polak, \emph{An Explicit Prime-Power Cone Criterion for the
Riemann Hypothesis}, version 1.2, 20 November 2026.
\href{https://zenodo.org/records/24406562}{Zenodo record 24406562};
\href{https://doi.org/10.5281/zenodo.24406562}{doi:10.5281/zenodo.24406562}.
\bibitem{Source}
L.~Alp\"oge and Claude, released Lean development accompanying
\href{https://alpo.ge/remark.pdf}{the announcement of 20 November 2026}.
\href{https://github.com/Deicyde/riemann-hypothesis/tree/40e56bab}{Source mirror, commit 40e56bab}.
Individual source files used in the extraction are linked in
Section \ref{sec:audit}.
\bibitem{Newman}
C.~M.~Newman, Fourier transforms with only real zeros,
\emph{Proceedings of the American Mathematical Society}
\textbf{61}(2) (1976), 245--251.
\href{https://sites.math.northwestern.edu/~auffing/papers/Newman.pdf}{Primary article PDF}.
\bibitem{PreviousNote}
\emph{Sharper sampling and waiting bounds for the prime-power cone criterion},
working note prepared for Robert Polak with ChatGPT,
21 November 2026, Theorem 6.1 (optimized asymptotic margin).
Unpublished derivation, not an external validation of the result.
\bibitem{MIT}
A.~Sutherland, \emph{Verification files for the Lean proof of the
Riemann Hypothesis}, MIT, update of 21 November 2026 at 08:45 UTC.
\href{https://math.mit.edu/~drew/grh/}{Primary verification report}.
\end{thebibliography}
\end{document}
